
\def\PUDcodigo{ECA208}
\def\PUDcodacad{04507.10}
\def\PUDdisciplina{Lógica de Programação}
\def\PUDhorasTeoria{40}
\def\PUDhorasPratica{40}

\def\PUDautor{Geraldo Luis Bezerra Ramalho}
\def\PUDdata{2013-04-17}
\def\PUDrevisao{4}
\def\PUDrevisor{Jonatha Rodrigues da Costa}
\def\PUDdatarevisao{2025-11-05}

\def\PUDsemestre{S2}
\def\PUDhoras{80}
\def\PUDcreditos{4}
\def\PUDprerequisito{ }%\hyperref[PUD:ACE003]{ACE003} }
\def\PUDpreacad{---}
%----------------------------------------------------------
\def\PUDementa{
	Estudo dos fundamentos de construção de algoritmos. Definição e uso de variáveis, tipos de dados, operadores e estruturas de controle (sequência, seleção e repetição). Tratamento de estruturas de dados homogêneas (vetores e matrizes). Introdução à modularização por meio de procedimentos e funções. Metodologias de desenvolvimento de programas e sua implementação em linguagem de alto nível, com ênfase nos processos de análise, teste e depuração de código. Conteúdo integrante da formação básica em Informática e Algoritmos.
}

%----------------------------------------------------------
\def\PUDconteudo{
	& Fundamentos de Algoritmos: Conceito de algoritmo; Técnicas de elaboração (pseudocódigo e fluxogramas); Tipos de dados, variáveis, constantes e operadores.
	\\
	& Estruturas de Controle: Estruturas sequenciais, condicionais e de repetição; Introdução à análise de complexidade básica de algoritmos.
	\\
	& Estruturas de Dados Elementares: Variáveis compostas homogêneas unidimensionais (vetores) e multidimensionais (matrizes).
	\\
	& Modularização e Funções: Procedimentos, funções e bibliotecas; Modularização para construção de programas estruturados.
	\\
	& Prática e Depuração: Introdução à Linguagem C; Implementação de algoritmos em linguagem de alto nível; Depuração de código e uso de ferramentas de compilação e análise.
}

%----------------------------------------------------------
\def\PUDhabilidades{
	\textbf{CG2 / HG21} – Modelar fenômenos e sistemas utilizando ferramentas computacionais.  
	Capacidade de transformar problemas de engenharia em modelos lógicos (algoritmos), utilizando sintaxe de programação para simulação de sistemas.
	
	\textbf{CG3 / HG31} – Conceber e projetar soluções viáveis tecnicamente.  
	Habilidade de conceber lógica estruturada para criação de soluções computacionais adequadas aos requisitos técnicos.
	
	\textbf{CG8 / HG81} – Atitude investigativa e autônoma.  
	Adoção de técnicas de depuração e otimização como práticas contínuas de investigação e melhoria do código.
}

%----------------------------------------------------------
\def\PUDatitudes{
	\textbf{Ética e Crítica (CG1/CG7):} Postura ética na manipulação e processamento de dados, reconhecendo a importância da integridade do \textit{software}. \\
	\textbf{Cooperação (CG6):} Comunicação clara e documentação adequada de algoritmos, favorecendo manutenção e trabalho em equipe. \\
	\textbf{Inovação (CG8):} Incentivo à criatividade e busca por soluções lógicas eficientes, com postura de atualização contínua.
}

%----------------------------------------------------------
\def\PUDmetodo{
	Aulas expositivas e práticas em laboratório. Serão utilizadas metodologias ativas, incluindo Aprendizagem Baseada em Problemas (ABP) e Desenvolvimento Baseado em Projetos, articulando teoria, prática e aplicação contextualizada.
}

%----------------------------------------------------------
\def\PUDrecursos{
	Projetor multimídia; Quadro branco e pincéis; Laboratório de Informática com computadores individuais para os estudantes; Acesso à internet; Ambiente de desenvolvimento e compiladores da linguagem C instalados.
	
}

\def\PUDavalia{
	A avaliação será contínua e formativa, composta por: provas teóricas, desenvolvimento e apresentação de programas funcionais (prova prática) e projetos em grupo. Instrumentos complementares poderão ser aplicados conforme necessidade.
}
%----------------------------------------------------------

\def\PUDobjetivos{
	Compreender os fundamentos de lógica de programação e desenvolvimento de programas estruturados.  Ler e escrever programas de computador utilizando linguagem de algoritmos e linguagem de alto nível. Além de verificar o contexto de aplicações regionais e interdiciplinariedade com diversas disciplinas presentes na matriz do curso.} 

\def\PUDbibbasica{
	 & \href{www.biblioteca.ifce.edu.br/index.asp?codigo_sophia=62925}{Ziviani}, N.  Projeto de algoritmos: com implementações em Java e C++.  São Paulo, SP : Cengage Learning, 2015.  621p.  ISBN: 9788522105250.
	\\ 
	& \href{www.biblioteca.ifce.edu.br/index.asp?codigo_sophia=24429}{Beneduzzi}, Humberto Martins.  Lógica e Linguagem de Programação: Introdução ao Desenvolvimento de Software.  144p.  ISBN: 9788563687111. 
	\\ 
	& \href{www.biblioteca.ifce.edu.br/index.asp?codigo_sophia=24502}{Souza}, M. A. F.; GOMES, M. M.; SOARES, M. V.; CONCILIO, R.  Algoritmos e Lógica de Programação.  São Paulo, SP : Cengage Learning, 2008.  214p.  ISBN: 8522104646. }

\def\PUDbibcomplementar{
	 & \href{www.biblioteca.ifce.edu.br/index.asp?codigo_sophia=24202}{Ziviani}, Nivio.  Projeto de algoritmos: com implementação em Pascal e C.  3º ed.  São Paulo, SP: Cengage Learning, 2011.  639p.  ISBN: 9788522110506.
	\\ 
	& \href{www.biblioteca.ifce.edu.br/index.asp?codigo_sophia=24181}{Schildt}, H.  C: completo e total.  3º ed.  São Paulo, SP : Pearson Makron Books, 1997.  827p.  ISBN: 9788534605953.
	\\ 
	&  \href{www.biblioteca.ifce.edu.br/index.asp?codigo_sophia=65446}{Cormen}, T. H.  Algoritmos: teoria e prática.  2º ed.  Rio de Janeiro, RJ: Elsevier, 2012.  926p.  ISBN: 9788535236996.
	\\
	&  \href{www.biblioteca.ifce.edu.br/index.asp?codigo_sophia=56907}{Ascencio}, Ana Fernanda Gomes; Campos, Edilene Aparecida Veneruchi de.  Fundamentos da Programação de Computadores.  E-book.  Pearson.  588p.  ISBN: 9788564574168.
	\\
	&  \href{www.biblioteca.ifce.edu.br/index.asp?codigo_sophia=40651}{Forbellone}, André Luiz Villar.  Lógica de programação: a construção de algoritmos e estruturas de dados.  3º ed.  São Paulo, SP: Pearson, 2013.  218p.  ISBN: 9788576050247. }


\PUDcorpo

%%--------------------------------------
% Anterior
%--------------------------------------
%\def\PUDementa{
%	Introdução ao conceito e desenvolvimento de algoritmos. Variáveis, tipos de dados, operadores, estruturas de controle (sequência, seleção, repetição). Vetores e Matrizes. Procedimentos e funções. Metodologias de desenvolvimento de programas e implementação em linguagem de alto nível, com ênfase na depuração de código. Conteúdo de Informática e Algoritmos (Conteúdo Básico).
%}
%
%\def\PUDconteudo{
%	& Fundamentos de Algoritmos: Conceito de algoritmo; Técnicas de elaboração (pseudocódigo e fluxogramas); Tipos de dados, variáveis, constantes e operadores.
%	\\
%	& Estruturas de Controle: Estruturas sequenciais, condicionais e de repetição; Introdução à análise de complexidade básica de algoritmos.
%	\\
%	& Estruturas de Dados Elementares: Variáveis compostas homogêneas unidimensionais (vetores) e multidimensionais (matrizes).
%	\\
%	& Modularização e Funções: Procedimentos, funções e bibliotecas; Modularização para construção de programas estruturados.
%	\\
%	& Prática e Depuração: Introdução à Linguagem C; Implementação de algoritmos em linguagem de alto nível; Depuração de código e uso de ferramentas de compilação e análise.
%}
%
%

%\def\PUDementa{\vm Introdução ao conceito de algoritmo. Desenvolvimento de algoritmos. Os conceitos de variáveis, tipos de dados, constantes, operadores aritméticos, expressões, atribuição, estruturas de controle (atribui- ção, seqüência, seleção, repetição). Vetores e Matrizes. Procedimentos, funções e bibliotecas. Metodologias de desenvolvimento de programas. Representação gráfica e textual de algoritmos. Estrutura e funcionalidades básicas de uma linguagem de programação procedimental. Implementação de algoritmos através da linguagem de programação introduzida. Depuração de Código e Ferramentas de Depuração.
%Introdução ao conceito de algoritmo.  Desenvolvimento de algoritmos.  Os conceitos de variáveis, tipos de dados, constantes, operadores aritméticos, expressões, atribuição, estruturas de controle (atribuição, seqüência, seleção, repetição).  Metodologias de desenvolvimento de programas.  Representação gráfica e textual de algoritmos.  Estrutura e funcionalidades básicas de uma linguagem de programação procedimental.  Implementação de algoritmos através da linguagem de programação introduzida.  Depuração de Código e Ferramentas de Depuração, Módulos (Procedimentos, Funções, Unidades ou Pacotes, Bibliotecas), Recursividade, Ponteiros e Alocação Dinâmica de Memória, Estruturas de Dados Heterogêneas (Registros ou Uniões, Arrays de Registros), Arquivos: Rotinas para manipulação de arquivos, Arquivos texto, Arquivos Binários, Arquivos de Registros.
%}

%\def\PUDconteudo{ & Técnicas de Elaboração de Algoritmos e Fluxogramas;  Algoritmos;   Fluxograma. 
%\\ 
% & Linguagem C;     Constantes: numérica, lógica e literal;             Variáveis: formação de identificadores, declaração de variáveis, comentários e comandos de atribuição;     Expressões e operadores aritméticos, lógicos, relacionais e literais, prioridade das operações;     Comandos de entrada e saída;     Estrutura seqüencial, condicional e de repetição. 
%\\ 
% & Estrutura de dados;     Variáveis compostas homogêneas unidimensionais (vetores).     Variáveis compostas homogêneas multidimensionais (matrizes).     %Variáveis compostas heterogêneas (registros).     Arquivos. 
%\\ 
% & Modularização;     Procedimentos e funções;    
% %Passagens de parâmetros;     Regras de escopo. 
% Aplicações com tecnologias atuais. 
%%\\ 
% %& Interfaces;     Porta paralela no PC;     Porta Serial RS232. 
% }
%\def\PUDmetodo{Aulas expositórias que podem ser teóricas e/ou práticas, onde as práticas no laboratório serão marcadas no decorrer da disciplina.}


%\def\PUDhabilidades{	
%	\textbf{CG2 / HG21} – Modelar fenômenos e sistemas utilizando ferramentas computacionais.  
%	Capacidade de transformar problemas de engenharia em modelos lógicos (algoritmos), utilizando sintaxe de programação para simulação de sistemas.
%	
%	\textbf{CG3 / HG31} – Conceber e projetar soluções viáveis tecnicamente.  
%	Habilidade de conceber lógica estruturada para criação de soluções computacionais adequadas aos requisitos técnicos.
%	
%	\textbf{CG8 / HG81} – Atitude investigativa e autônoma.  
%	Adoção de técnicas de depuração e otimização como práticas contínuas de investigação e melhoria do código.
%}


%\def\PUDavalia{A avaliação será feita com aplicação de provas teóricas e/ou práticas, além da possibilidade de inclusão de trabalhos e seminários no decorrer da disciplina.}

%%--------------------------------------